\subsection{The weight-lowering Bockstein}

Let $R=\mathbb Z_{(p)}$, and let $M_w(R)$ denote the space of level-one
modular forms of weight $w$ with coefficients in $R$.  We identify a form
with its $q$-expansion.  Write
\[
  \overline M_w:=M_w(R)/pM_w(R)
\]
and let $\overline{\mathcal D}$ be the space of divided congruences modulo
$p$, viewed through their $q$-expansions.

\begin{definition}\label{def:bockstein}
Let $g\in\mathcal D(R/p^2R)$ and suppose that $g\pmod p$ is represented by
some $h\in M_w(R)$.  Define
\[
  \beta_w(g):=\left[\frac{g-h}{p}\right]
       \in \overline{\mathcal D}/\overline M_w.
\]
Here any lift of the $q$-expansion of $g$ to $R$ may be used in the
numerator.
\end{definition}

\begin{lemma}\label{lem:bockstein}
The class $\beta_w(g)$ is independent of the choices in Definition
\ref{def:bockstein}.  Moreover,
\[
  \beta_w(g)=0
  \quad\Longleftrightarrow\quad
  g\equiv H\pmod{p^2}
  \quad\text{for some }H\in M_w(R).
\]
Consequently, if $w$ lies in the admissible weight class of $g$, then
\[
  \beta_w(g)=0
  \quad\Longleftrightarrow\quad
  \omega_{p^2}(g)\le w.
\]
\end{lemma}

\begin{proof}
Suppose first that $h,h'\in M_w(R)$ both represent $g\pmod p$.  Then
$h-h'$ has $q$-expansion divisible by $p$.  The $q$-expansion lattice
$M_w(R)$ is saturated: an integral modular form whose coefficients are all
divisible by $p$ belongs to $pM_w(R)$.  One can also see this directly by
forward substitution in the standard integral echelon basis, whose pivot
coefficients are units.  Thus $h-h'\in pM_w(R)$.  Hence
\[
  \frac{g-h}{p}-\frac{g-h'}{p}=\frac{h'-h}{p}\in M_w(R),
\]
so the two choices define the same class modulo $\overline M_w$.  Changing
the integral lift of $g$ changes the quotient by an arbitrary multiple of
$p^2/p=p$, and hence does not change its reduction modulo $p$.  This proves
well-definedness.

If $\beta_w(g)=0$, there is an $m\in M_w(R)$ such that
$(g-h)/p\equiv m\pmod p$.  Thus
\[
  g\equiv h+pm\pmod{p^2},
\]
and $h+pm$ belongs to $M_w(R)$.  Conversely, suppose that
$g\equiv H\pmod{p^2}$ for some $H\in M_w(R)$.  Since
$H\equiv h\pmod p$, the same echelon-basis argument gives
$H-h=pm$ for some $m\in M_w(R)$.  It follows that
\[
  \frac{g-h}{p}\equiv\frac{H-h}{p}=m\pmod p,
\]
so $\beta_w(g)=0$ in
$\overline{\mathcal D}/\overline M_w$.  For the final equivalence, assume
that every smaller admissible representative weight differs from $w$ by a
nonnegative multiple of $P=p(p-1)$.  Multiplication by the corresponding
power of $E_{p-1}^p\equiv1\pmod {p^2}$ raises such a representative to
weight $w$ without changing its reduction modulo $p^2$.  Therefore a
representative of weight at most $w$ exists exactly when one of weight $w$
exists.  This is the asserted filtration criterion.
\end{proof}

\begin{remark}\label{rem:bockstein-computation}
The canonical echelon remainder used in the computations is a coordinate
representative of $\beta_w(g)$.  Reducing $g$ modulo the integral echelon
basis of $M_w(R)$ kills the same leading coefficients as the lift $h$ in
the proof.  The resulting remainder is divisible by $p$, and division by
$p$ followed by reduction modulo $p$ gives the displayed class.  Thus the
computational zero test is intrinsic even though its coordinate vector
depends on the chosen echelon basis.
\end{remark}
